% Source PDF: source/普通高考/2016/2016天津理.pdf, page 1.
% PDF embedded bitmap attempt: pdfimages -list returned no image objects; the flowchart is vector art.
% Grid evidence: tmp/redraw_flowchart_2016_tianjin_science/grid/page1_grid100.png and q4_block_grid50.png.
% Comparison source: tmp/redraw_flowchart_2016_tianjin_science/crops/q4_orig_final.png,
% cropped from the ungridded page render with a safety border.
% Elements: rounded start/end terminals; S=4 and n=1 initialization boxes; two
% decision diamonds S>=6? and n>3?; branch updates S=2S and S=S-6; n=n+1;
% output parallelogram; yes/no labels and loop-back connectors.
% Relations: start -> S=4 -> n=1 -> S>=6?; no -> S=2S, yes -> S=S-6;
% both update branches merge into n=n+1 -> n>3?; no loops to S>=6?, yes -> output S -> end.
% Directed edges: start->S=4, S=4->n=1, n=1->S>=6; S>=6-(否)->S=2S;
% S>=6-(是)->S=S-6; both updates->n=n+1; n=n+1->n>3;
% n>3-(否)->S>=6; n>3-(是)->output; output->end.
% solid segments: all node borders and directed connectors; dashed segments: none.
% labels: formulas centered in nodes; branch labels sit beside the corresponding
% diamond edges without hand-spacing or overlap.
\begin{tikzpicture}[x=1cm,y=1cm,baseline,
  flowarrow/.style={-{Latex[length=2.1mm,width=1.5mm]},line width=0.45pt},
  nodebox/.style={draw,line width=0.45pt,inner sep=2pt,minimum height=0.62cm,anchor=center,align=center,
    execute at begin node={\setlength{\parindent}{0pt}\noindent}},
  branchlabel/.style={anchor=center,align=center,inner sep=0pt,
    execute at begin node={\setlength{\parindent}{0pt}\noindent}},
  term/.style={nodebox,rounded corners=2.5pt,minimum width=1.15cm,text width=1.15cm},
  io/.style={nodebox,shape=trapezium,trapezium left angle=70,trapezium right angle=110,
    minimum width=1.95cm,text width=1.65cm,minimum height=0.72cm},
  decision/.style={nodebox,diamond,aspect=2.9,minimum width=3.15cm,minimum height=0.92cm,inner sep=1pt}
]
  % Main top-to-bottom program flow.
  \node[term] (start) at (0,0) {开始};
  \node[nodebox,text width=1.65cm,minimum width=1.65cm] (sinit) at (0,-1.20) {$S=4$};
  \node[nodebox,text width=1.65cm,minimum width=1.65cm] (ninit) at (0,-2.40) {$n=1$};
  \node[decision] (testS) at (0,-3.70) {$S\geq 6?$};
  \node[nodebox,text width=2.00cm,minimum width=2.00cm] (doubleS) at (-2.55,-5.15) {$S=2S$};
  \node[nodebox,text width=2.35cm,minimum width=2.35cm] (subtract) at (0,-5.15) {$S=S-6$};
  \node[nodebox,text width=2.35cm,minimum width=2.35cm] (inc) at (0,-6.60) {$n=n+1$};
  \node[decision] (testN) at (0,-8.00) {$n>3?$};
  \node[io] (output) at (0,-9.35) {输出 $S$};
  \node[term] (finish) at (0,-10.55) {结束};

  \draw[flowarrow] (start.south) -- (sinit.north);
  \draw[flowarrow] (sinit.south) -- (ninit.north);
  \draw[flowarrow] (ninit.south) -- (testS.north);
  \draw[flowarrow] (testS.south) -- (subtract.north);
  \draw[flowarrow] (inc.south) -- (testN.north);
  \draw[flowarrow] (testN.south) -- (output.north);
  \draw[flowarrow] (output.south) -- (finish.north);

  % 否 branch to S=2S, then merge with the 是 branch before n=n+1.
  \draw (testS.west) -- (-2.55,-3.70);
  \draw[flowarrow] (-2.55,-3.70) -- (doubleS.north);
  % Keep the branch arrow at the three-way merge; the merged path has its own
  % arrow at the actual increment-node endpoint below.
  \draw[flowarrow] (doubleS.south) -- (-2.55,-5.85) -- (-0.45,-5.85) -- (0,-5.85);
  \draw (subtract.south) -- (0,-5.85);
  \draw[flowarrow] (0,-5.85) -- (inc.north);

  % n>3 否 returns to the first decision; the output route stays independent.
  \draw[flowarrow] (testN.east) -- (3.05,-8.00) -- (3.05,-2.85) -- (0,-2.85);

  \node[branchlabel] at (-2.22,-3.12) {否};
  \node[branchlabel] at (0.42,-4.28) {是};
  \node[branchlabel] at (2.15,-7.54) {否};
  \node[branchlabel] at (0.42,-8.72) {是};
\end{tikzpicture}
